%1/2in left/right margins, 1in top, 1/2 bottom
\documentclass[12pt]{article}
\hbadness=10000
\tolerance=10000
\textheight 9.7in
\textwidth 7.0in
\oddsidemargin -.5in
\topmargin -.5in
\headsep 0in
\headheight 0in
%\pagestyle{empty}

\newcommand{\by}{\mbox{\boldmath$y$}}
\newcommand{\bX}{\mbox{\boldmath$X$}}
\newcommand{\bbeta}{\mbox{\boldmath$\beta$}}

\begin{document}
\hrule
\medskip
\centerline{\textbf{STATISTICS 801 - Section 250}}\hfill \textbf{Fall 2016}
\centerline{\textbf{z Test for a Proportion}}
\smallskip
\hrule
Many hypothesis testing situations involve proportions. A hypothesis test involving a population proportion can be considered as a binomial experiment when there are only two outcomes and the probability of success does not change from trial to trial.\\
For a binomial random variable $E(X)=np$ and $V(X)=np(1-p)$.\\
Since the normal distribution can be used to approximate the binomial distribution when $np\geq 5$ and $n(1-p)\geq 5$, the standard normal distribution can be used to test hypotheses for proportions.\\
Test statistic for proportions:\\
\begin{equation}
z=\frac{\hat{p}-p}{\sqrt{p(1-p)/n}}
\end{equation}
where $\hat{p}=\frac{X}{n}$ (sample proportion),\\
p is the population proportion, and n is the sample size.

\begin{enumerate}
\item An educator estimates that the dropout rate for seniors at high schools in Nebraska is $15\%$. Last year, 38 seniors from a random sample of 200 Nebraska seniors withdrew. At $\alpha=0.05$, is there enough evidence to reject the educator's claim?
\begin{enumerate}
\item State the hypotheses.\\
$H_0:p=0.15$ and $H_1:p\neq 0.15$
\item Check whether we can use the standard normal distribution to test the hypothesis.\\
$np=38\times 0.15=5.7$ and $n(1-p)=38\times 0.85=32.3$
\item Find the critical value(s). Since $\alpha=0.05$ and the test is two-tailed, the critical values are 1.96 and -1.96.
\item Compute the test statistic. First, it is necessary to find $\hat{p}$.\\
$\hat{p}=\frac{X}{n}=\frac{38}{200}=0.19$ and $p=0.15$, so $1-p=0.85$\\
$z=\frac{\hat{p}-p}{\sqrt{p(1-p)/n}}=\frac{0.19-0.15}{\sqrt{(0.15)(0.85)/200}}=1.58$
\item Make the decision. We fail to reject the null hypothesis because the test statistic falls outside of the rejection region.
\item Make a conclusion. There is not enough evidence to reject the claim that the dropout rate for seniors in high schools in Nebraska is $15\%$.\\
\\
\\
\\
\\
\\
\\
\\
\\
\\
\\
\\
\\
\end{enumerate}
\item An attorney claims that more than $25\%$ of all lawyers advertise. A sample of 200 lawyers in a certain city showed that 63 had used some form of advertising, At $\alpha=0.05$, is there enough evidence to support the attorney's claim? Use the p-value method.
\begin{enumerate}
\item State the hypotheses.\\
$H_0: \;\;\;\;\;\;\;\;\;\;\;\;\;\;\;\;\;\;\;\;\;\;\;\;\;\; $ and $H_1:$
\item $\hat{p}=$\\
\\
$p= \;\;\;\;\;\;\;\;\;\;\;\;\;\;\;\;\;\;\;\;\; 1-p=$\\
$z=$\\
\\
\\
\item Find the p-value.\\
\\
\item Make a decision.\\
\\
\item Make a conclusion.\\
\\
\\
\\
\end{enumerate}
A recent survey found that $64.7\%$ of the population own their homes. In a random sample of 150 heads of households, 92 responded that they owned their homes. At the 0.01 level of significance, does that indicate a difference from the national proportion?
\end{enumerate}
\end{document}

